\section*{Overview}

Digit DP is the standard way to count or optimize over huge numeric ranges when the property depends on the decimal
digits themselves. The main shift is to stop thinking about numbers as integers and start thinking about them as
left-to-right digit constructions under an upper-bound constraint.

\section*{Problem-Driven Motivation}

Suppose \(R\) can be as large as \(10^{18}\), and we must count numbers in \([L,R]\) whose digits satisfy some rule:

\begin{itemize}
  \item no adjacent equal digits,
  \item digit sum equals \(S\),
  \item remainder modulo \(m\) is fixed,
  \item the decimal expansion avoids some forbidden pattern.
\end{itemize}

Brute force is impossible. The only thing that makes the problem manageable is that while building a number from left to
right, the future only depends on a small amount of state about the prefix.

\section*{Recognition Pattern}

Digit DP is the right tool when:

\begin{itemize}
  \item the numeric bound is enormous,
  \item the validity condition is defined on digits,
  \item the natural range answer is \(\texttt{solve}(R) - \texttt{solve}(L-1)\),
  \item the important memory about the prefix is small enough to memoize.
\end{itemize}

If the property is not about digits at all, forcing a digit-DP model usually makes the problem worse.

\section*{Derivation}

Write the upper bound \(N\) as a digit string. Process it from most significant digit to least significant digit.

The standard state has three conceptual parts:

\begin{itemize}
  \item \textbf{position}: which digit are we choosing,
  \item \textbf{tight}: whether the built prefix is still exactly equal to the prefix of \(N\),
  \item \textbf{problem state}: whatever information about the chosen digits still matters.
\end{itemize}

The \texttt{tight} flag is what enforces the upper bound. While it is true, the next digit cannot exceed the
corresponding digit of \(N\). Once it becomes false, the remaining digits are unconstrained.

Leading zeros need special care. In many problems, ``we have not started the number yet'' must be represented
explicitly, or folded into a sentinel previous digit.

\section*{Worked Problem}

\paragraph{Problem.}
Count integers in \([L,R]\) whose decimal representation has no two adjacent equal digits.

\paragraph{Why naive fails.}
The interval can be enormous, so scanning every number is impossible.

\paragraph{State design.}
Process digits left to right and keep:

\begin{itemize}
  \item \texttt{pos},
  \item \texttt{tight},
  \item \texttt{prev} = previous chosen digit, or a sentinel \(10\) meaning ``no real digit has started yet''.
\end{itemize}

\paragraph{Transition.}
At each position choose the next digit \(d\) between \(0\) and the current limit. The transition is illegal only if:

\begin{itemize}
  \item the number has already started,
  \item and \(d = prev\).
\end{itemize}

If the number has not started and we place another leading zero, the sentinel \(10\) stays.

\paragraph{Why memoization works.}
When \texttt{tight = false}, the future depends only on \((pos, prev)\). Two states with the same values have identical
continuations, so they can share one memoized answer.

\section*{Implementation Reasoning}

The important implementation habit is to write \(\texttt{solve}(N)\) first and keep range logic separate:
\[
answer(L,R) = solve(R) - solve(L-1).
\]

For the no-adjacent-equal example, the sentinel previous digit removes the need for a separate \texttt{started} flag.
That is often cleaner than carrying two booleans.

The code below is a reusable digit-DP skeleton for this style of problem:

\begin{itemize}
  \item digits are processed left to right,
  \item memoization is used only for non-tight states,
  \item one extra state parameter (\texttt{prev}) captures the property.
\end{itemize}

\section*{Correctness Intuition}

Every valid number in \([0,N]\) corresponds to exactly one root-to-leaf path in the digit DP. The \texttt{tight} flag
ensures that invalid paths exceeding \(N\) are never explored. The problem state stores exactly the information needed
to judge future validity, so memoized non-tight states are genuinely identical subproblems.

\section*{Complexity Analysis}

If the custom state space has size \(S\), the rough bound is
\[
O(\text{digits} \cdot S \cdot 10).
\]

For the no-adjacent-equal example, \(S\) is tiny: 11 possible previous-digit states.

\section*{Common Pitfalls}

\begin{itemize}
  \item Memoizing tight states even though they depend on the exact upper-bound prefix.
  \item Mishandling leading zeros and accidentally treating them as real digits.
  \item Forgetting the \(\texttt{solve}(R) - \texttt{solve}(L-1)\) pattern.
  \item Letting the custom state explode until the DP is no longer small.
\end{itemize}

\section*{Variants and Failure Modes}

\begin{itemize}
  \item Replace \texttt{prev} by digit sum, remainder modulo \(m\), bitmask of used digits, or automaton state.
  \item More complicated properties may need both \texttt{started} and a custom state.
  \item If the property depends on the number in a non-digit way, digit DP is the wrong abstraction.
\end{itemize}

\section*{Practice Problems}

\begin{itemize}
  \item Count numbers with no adjacent equal digits.
  \item Count numbers with digit sum or remainder restrictions.
  \item Count numbers avoiding a forbidden decimal substring.
\end{itemize}

\section*{References}

\begin{itemize}
  \item \href{https://usaco.guide/gold/digit-dp?lang=cpp}{USACO Guide: Digit DP}.
  \item \href{https://codeforces.com/catalog}{Codeforces Catalog}.
  \item \href{https://oiwiki.org/dp/number/}{OI Wiki: Digit DP}.
\end{itemize}
